\documentclass[12pt,a4paper]{article}

% --- Paquetes de Idioma y Codificación ---
\usepackage[utf8]{inputenc}
\usepackage[T1]{fontenc}
\usepackage[spanish,es-tabla]{babel}
\usepackage{csquotes}

% --- Paquetes de Estética y Formato ---
\usepackage[margin=2.54cm]{geometry}
\usepackage{times}
\usepackage{setspace}
\usepackage{titlesec}
\usepackage{enumitem}
\usepackage{graphicx}
\graphicspath{{./}{figuras/}{figuras/figuras1/}{figuras/figuras3/}{figuras/figuras4/}{figuras/figuras5/}{figuras/figuras6/}}
\usepackage{tikz}
\usetikzlibrary{shapes.geometric, arrows.meta, positioning, fit, backgrounds, calc, trees}
\usepackage{caption}
\usepackage{multirow}
\usepackage{pdflscape}
\usepackage[table]{xcolor}
\usepackage{fontawesome5}
\usepackage{float}
\usepackage{longtable}
\usepackage{booktabs}
\usepackage{array}
\usepackage{indentfirst}
\usepackage{amsmath} % NECESARIO para los comandos \text en las hipótesis
\usepackage{tabularx}
% Dependencias del capítulo 2 (sin alterar el formato del capítulo 1).
\usepackage{listings}
\usepackage{needspace}
\usepackage{ragged2e}
\newcolumntype{Y}{>{\RaggedRight\arraybackslash}X}
\newcolumntype{P}[1]{>{\RaggedRight\arraybackslash}p{#1}}

\usepackage{amsmath}
\usepackage{amssymb}
\usepackage{enumitem}

% --- Configuración de Párrafo ---
\onehalfspacing 
\setlength{\parindent}{1.27cm}
\setlength{\parskip}{0pt}

% --- Títulos Estilo APA 7 ---
\titleformat{\section}{\normalfont\large\bfseries\centering}{\thesection}{1em}{}
\titleformat{\subsection}{\normalfont\normalsize\bfseries}{\thesubsection}{1em}{}
\titleformat{\subsubsection}{\normalfont\normalsize\bfseries\itshape}{\thesubsubsection}{1em}{}

% --- Configuración de Bibliografía ---
\usepackage[backend=biber, style=apa]{biblatex}
\addbibresource{referencias.bib}
\addbibresource{practicas/practica7_fuentes.bib}
\addbibresource{practicas/practica8_fuentes.bib}

% --- Enlaces ---
\usepackage{hyperref}
\hypersetup{
    colorlinks=true,
    allcolors=black, 
    citecolor=black
}

% --- Numeración de Páginas ---
\usepackage{fancyhdr}
\setlength{\headheight}{14.5pt}
\pagestyle{fancy}
\fancyhf{}
\renewcommand{\headrulewidth}{0pt}
\fancyhead[R]{\thepage}

% --- Comando personalizado para Páginas Divisoras de Capítulos (CORREGIDO) ---
\newcommand{\capitulopage}[1]{%
    \newpage
    \thispagestyle{empty} % Quita la numeración de página en esta hoja
    \vspace*{\fill}
    \begin{center}
        {\normalfont\Large\bfseries #1 \par}
    \end{center} % <--- Aquí estaba el error, ahora cierra correctamente
    \vspace*{\fill}
    \newpage
}

% --- Configuración de formato para la numeración de títulos ---
\renewcommand{\thesection}{\Roman{section}} % Secciones en Romanos (I, II, III)
\renewcommand{\thesubsection}{\arabic{section}.\arabic{subsection}} % Subsecciones en Arábigos (1.1, 1.2)
\renewcommand{\thesubsubsection}{\arabic{section}.\arabic{subsection}.\arabic{subsubsection}} % Sub-subsecciones (1.1.1)

\begin{document}

% --- CUBIERTA / DATOS INFORMATIVOS ---
\begin{titlepage}
    \centering
    \setstretch{1.5}

    \includegraphics[width=0.3\textwidth]{lognuevo_ucsm.png}

    % --- ENCABEZADO ---
    \vspace{0.5cm}
    {\large\bfseries FACULTAD DE CIENCIAS E INGENIERÍAS FÍSICAS Y FORMALES \par}
    {\large\bfseries ESCUELA PROFESIONAL DE INGENIERÍA DE SISTEMAS \par}
    \vspace{0.4cm}
    
    {\large Análisis estático de alcanzabilidad mediante grafos de llamadas para la reducción de alertas de vulnerabilidad en dependencias transitivas de Python - Lidercom Perú S.A.C.\par}
        
    \vspace{0.5cm}
    
    % --- PRESENTADO POR ---
    {\large\bfseries Presentado por: \par}
    {\large Centeno León, Julio Enrique \par}
    {\large Yáñez Mejía, José Alonso \par}    
    \vspace{0.5cm}
    
    % --- ASESOR ---
    {\large\bfseries Asesor: \par}
    {\large Dr. Néstor Torres Gamarra \par}
    
    \vfill

    % --- LUGAR Y FECHA ---
    {\large Arequipa -- Perú \par}
    {\large 2026 \par}
\end{titlepage}

\newpage
\tableofcontents
\newpage

\include{capitulos/capitulo1}
\include{capitulos/capitulo2}


\newpage
\section{REFERENCIAS}
\printbibliography[heading=none]

% =========================================================================
% APÉNDICE A. GLOSARIO DE TÉRMINOS DEL DOMINIO
% =========================================================================
\newpage
\appendix
\section*{Apéndice A}
\addcontentsline{toc}{section}{Apéndice A}

\begin{center}
    {\normalfont\Large\bfseries Apéndice A. Glosario de Términos del Dominio \par}
\end{center}
\vspace{0.3cm}

\noindent
En cumplimiento de las Unidades 3 y 4 del libro de cátedra y las normas de presentación APA 7.ª edición \parencite{apa2020publication, torres2026proyecto}, a continuación se definen los términos propios del dominio de la ciberseguridad, análisis estático y cadena de suministro de software aplicados en el proyecto \textit{PyReach}, divididos en dos secciones continuas para optimizar su legibilidad.

% --- PARTE 1 ---
\begin{table}[H]
\centering
\caption{Glosario de términos del dominio (Parte 1)}
\label{tab:glosario_terminos_u3_u4}
\small
\setstretch{1.05}
\begin{tabularx}{\textwidth}{>{\bfseries\raggedright\arraybackslash}p{4.4cm} X}
\toprule
\textbf{Término} & \textbf{Definición} \\
\midrule
Alerta No Explotable & 
Notificación de seguridad emitida por un escáner SCA tradicional sobre una vulnerabilidad presente en una librería que no puede ser activada por un atacante, debido a que el código vulnerable no es alcanzable desde la aplicación huésped \parencite{pashchenko2022vuln4real}. \\
\addlinespace
Análisis de Alcanzabilidad (\textit{Reachability Analysis}) & 
Técnica de análisis formal de programas que determina estáticamente si existe una ruta de flujo de control o camino dirigido válido que conecte los puntos de entrada (\textit{entrypoints}) de la aplicación con una función vulnerable específica \parencite{dann2021identifying, ponta2020detection}. \\
\addlinespace
Árbol de Sintaxis Abstracta (AST) & 
Representación jerárquica en forma de árbol de la estructura sintáctica y gramatical del código fuente interpretado, generada por el analizador sintáctico del lenguaje sin incluir delimitadores ni espacios en blanco \parencite{salis2021pycg}. \\
\addlinespace
Compuerta de Calidad (\textit{Quality Gate}) & 
Punto de control automatizado en un pipeline de integración continua que evalúa métricas o resultados de pruebas estáticas y detiene el pase a producción mediante códigos de salida (\textit{exit codes}) si no se cumplen los criterios mínimos definidos \parencite{humble2010continuous}. \\
\addlinespace
Dependencia Transitiva & 
Paquete o módulo de software requerido de forma indirecta por una aplicación, al ser una dependencia secundaria declarada e importada por otra librería de primer orden (\textit{dependencia directa}) \parencite{decan2019empirical, alfadel2023empirical}. \\
\addlinespace
Ejecución \textit{On-Premise} & 
Modelo de procesamiento y despliegue donde el software y los análisis se ejecutan estrictamente en servidores locales o infraestructura interna de la organización, garantizando que cero bytes de datos abandonen la red corporativa \parencite{ladisa2023sok, gutierrez2024auditoria}. \\
\bottomrule
\end{tabularx}
\vspace{2pt}
\begin{flushleft}
\scriptsize \textit{Nota.} Términos del dominio estructurados según la Tabla 26 del modelo institucional y las Unidades 3 y 4 del libro de cátedra \parencite{torres2026proyecto}. Elaboración propia.
\end{flushleft}
\end{table}

\newpage

% --- PARTE 2 ---
\begin{table}[H]
\centering
\ContinuedFloat
\caption{Glosario de términos del dominio (Continuación)}
\small
\setstretch{1.05}
\begin{tabularx}{\textwidth}{>{\bfseries\raggedright\arraybackslash}p{4.4cm} X}
\toprule
\textbf{Término} & \textbf{Definición} \\
\midrule
Fatiga de Alertas (\textit{Alert Fatigue}) & 
Fenómeno psicológico y operativo que experimentan los ingenieros de software ante un volumen masivo y continuo de alertas irrelevantes o falsos positivos, provocando que desatiendan o ignoren incidentes de seguridad críticos reales \parencite{johnson2013developers, sadowski2018lessons}. \\
\addlinespace
Grafo de Llamadas (\textit{Call Graph}) & 
Grafo dirigido cuyos nodos representan funciones, métodos o procedimientos, y cuyas aristas dirigidas modelan las relaciones de invocación y transferencia de control de ejecución entre ellos \parencite{ryder1979constructing, salis2021pycg}. \\
\addlinespace
OpenSSF OSV & 
Estándar abierto y base de datos distribuida (\textit{Open Source Vulnerabilities}) impulsada por la Open Source Security Foundation, que documenta avisos de seguridad estructurados en JSON con rangos de versiones y símbolos afectados \parencite{openssf2023osv}. \\
\addlinespace
Punto de Entrada (\textit{Entrypoint}) & 
Función, script principal o manejador de ruta (ej. endpoints de FastAPI o vistas de Django) donde se inicia la ejecución del código de la aplicación tras recibir una petición externa o invocación del usuario \parencite{salis2021pycg}. \\
\addlinespace
SARIF (\textit{Static Analysis Results Interchange Format}) & 
Estándar internacional de OASIS basado en JSON para la representación estructurada e interoperable de los diagnósticos y trazas emitidos por herramientas de análisis estático en pipelines de CI/CD \parencite{oasis2019sarif}. \\
\addlinespace
Triaje Manual & 
Proceso de revisión humana en el que un desarrollador o especialista en ciberseguridad inspecciona, verifica y clasifica individualmente una alerta de vulnerabilidad para dictaminar si representa un riesgo real o un falso positivo \parencite{pashchenko2022vuln4real, sadowski2018lessons}. \\
\bottomrule
\end{tabularx}
\vspace{2pt}
\begin{flushleft}
\scriptsize \textit{Nota.} Continuación de la Tabla 26 conforme a los lineamientos metodológicos de cátedra \parencite{torres2026proyecto}. Elaboración propia.
\end{flushleft}

\end{table}

\end{document}